% Card: cutting off the king (rook endings)
\paragraph{Cut off}\index{Rook endings!Cut-off}

\diagram[cut-off]{4k3/1R2P3/r7/3K4/8/8/8/8 w - - 1 5}

White's pawn is one step from queening, but the black king stands in front of it, and the black
rook on the sixth rank keeps the white king away. The king would have to cross the sixth rank to
come closer to its pawn, and it cannot. The black rook has
\vocab[rook endings]{cut-off}{cut off}{A rook that controls a whole rank or file cuts off the
enemy king: the king cannot cross that line. The defender uses it to keep the attacking king away
from its pawn, the attacker to keep the defending king away from the pawn's path}
the white king.

White cannot make progress, and the position is a draw. As long as the rook keeps guarding the
sixth rank, the white king cannot get closer than the fifth rank.
